\section{Multi-cell operations}

The statements described in the previous section update at most one cell at a
time, and did not allow writing to a cell that already holds a value.
However, sometimes we need to go beyond this.  The most common case is an
operation to sort a block of cells, according to some criterion.  

The language contains a built in function with signature
\begin{scala}
def sortBlockBy(columns: List[Column], rows: List[Row], 
                before: (Row,Row) => Boolean): Unit
\end{scala}
%
This sorts the block defined by |columns| and |rows|, according to the
criterion defined by |before|: if |before(r1,r2)| gives true, then row~|r1|
will be put before~|r2| in the sorted version (so |before| should correspond
to an order relation). 

Figure~\ref{fig:sortByTotal} illustrates this function.  The script expects to
find names in column |#A| from row |#1|, with corresponding |Int| scores in
columns |#B| to |#D|.  It calculates the total for each row, putting the
results in column |#E|.  However, we want ways to sort the data according to
the total or by the name. 


% #include "haskell.dir"
\begin{figure}
\begin{scala}
#include "spreadsheet.dir"
val firstEmpty = firstEmptyInColumn(#A)  // Find end of inputs.
val candidateRows = #1 until firstEmpty // Rows for candidates.
val columns = #A to #E  // Columns in the table.
// Calculate totals.
for(r <- candidateRows) Cell(#E,r) = Cell(#B,r):Int + Cell(#C,r) + Cell(#D,r)
// Sorting operations.
operation sortByTotal() = {
  def before(r1: Row, r2: Row) = Cell(#E,r1): Int >= Cell(#E, r2)
  sortBlockBy(columns, candidateRows, before)
}
operation sortByName() = {
  def before(r1: Row, r2: Row) = Cell(#A,r1): String <= Cell(#A, r2)
  sortBlockBy(columns, candidateRows, before)
}
\end{scala}
\caption{A script illustrating the {\scalashape sortByTotal} function.}
\label{fig:sortByTotal}
\end{figure}

The operation |sortByTotal()| provides a way to sort the data into decreasing
order of total score (column~|#E|): the |before| function defines that
row~|r1| will be placed before row~|r2| if it has a higher total score.  

If we had placed the body of this operation at the top level of the script,
i.e.~not in any function or operation, then it would have been executed
\emph{every} time that the user entered a value in a cell.  This would have
been distracting for the user, because the row they were currently typing into
could be moved.

Instead, the keyword |operation| means that an item labelled ``|sortByTotal|''
is added to the ``Operations'' menu in the application.  Selecting this item
invokes the operation, and so sorts the data.

\framebox{screenshot}

The operation |sortByName| is similar, but sorts the data into decreasing
order by the name field.
